一行公式可以装下半天课的内容。

比如梯度下降的更新式：\(\theta_{t+1}=\theta_t-\eta\nabla L(\theta_t)\)。背下来不过十秒。但你拿着这一行公式面对实际问题时，它会沉默——变量属于什么空间？损失函数为什么要写成这个样子？导数在哪儿存在、在哪儿断裂？浮点计算会不会把梯度吞成零？训练数据和部署数据之间那层薄薄的假设，撑得住吗？

一行公式压缩了所有这些东西，却一个也不显示。

再看一个很小的例子。用直线 \(y=\beta_0+\beta_1x\) 拟合三个点
\[
(0,1),\qquad (1,2),\qquad (2,2).
\]
把每个预测写成``截距加斜率乘输入''，就得到
\[
X=
\begin{bmatrix}
1&0\\
1&1\\
1&2
\end{bmatrix},
\qquad
y=
\begin{bmatrix}
1\\2\\2
\end{bmatrix}.
\]
第一列全为 \(1\)，负责截距；第二列记录输入 \(x\)，负责斜率。最小二乘不是要求直线穿过三个点——这通常做不到——而是让三个竖直残差的平方和最小。

设参数向量为 \(\beta=(\beta_0,\beta_1)^{\trans}\)。三个预测值可以一次写成 \(X\beta\)，所以残差向量和要最小化的平方和分别是
\[
r(\beta)=y-X\beta,
\qquad
S(\beta)=\norm{y-X\beta}^{2}
=\sum_{i=1}^{3}(y_i-x_i^{\trans}\beta)^2,
\]
其中 \(x_i^{\trans}\) 是 \(X\) 的第 \(i\) 行。现在追问：为什么最优解会出现 \(X^{\trans}X\)？

先对任意一个参数 \(\beta_j\) 求偏导。第 \(i\) 个残差对 \(\beta_j\) 的变化率是 \(-X_{ij}\)，由链式法则
\[
\frac{\partial S}{\partial\beta_j}
=-2\sum_{i=1}^{3}X_{ij}(y_i-x_i^{\trans}\beta).
\]
把 \(j=0,1\) 的两个偏导叠成一个向量，求和恰好就是一次矩阵转置乘法：
\[
\nabla_\beta S(\beta)
=-2X^{\trans}(y-X\beta).
\]
有些教材把转置写成 \(X'\)，所以 \(X'X\) 与这里的 \(X^{\trans}X\) 是同一个对象。转置并非凭空加入：\(X\) 先把参数变化变成三个预测值的变化，\(X^{\trans}\) 再把三个残差分别沿截距和斜率两个参数方向汇总回来。

最小点处若没有约束，梯度必须为零。因此
\[
\begin{aligned}
-2X^{\trans}(y-X\hat\beta)&=0,\\
X^{\trans}y-X^{\trans}X\hat\beta&=0,\\
X^{\trans}X\hat\beta&=X^{\trans}y.
\end{aligned}
\]
最后一行就是正规方程。这里 \(X^{\trans}X\) 出现的原因很具体：把括号中的 \(X\hat\beta\) 左乘 \(X^{\trans}\)，结合律给出 \(X^{\trans}(X\hat\beta)=(X^{\trans}X)\hat\beta\)。

代入本例的数据，
\[
X^{\trans}X=
\begin{bmatrix}
3&3\\
3&5
\end{bmatrix},
\qquad
X^{\trans}y=
\begin{bmatrix}
5\\6
\end{bmatrix},
\]
若 \(X\) 的两列线性无关，则 \(X^{\trans}X\) 可逆，因此
\[
\hat\beta=(X^{\trans}X)^{-1}X^{\trans}y
=
\begin{bmatrix}
7/6\\[2pt]1/2
\end{bmatrix}.
\]
拟合直线是 \(\hat y=7/6+x/2\)。三个预测值为 \(7/6,5/3,13/6\)，残差为 \((-1/6,1/3,-1/6)^{\trans}\)；它们并不全为零，但满足 \(X^{\trans}(y-X\hat\beta)=0\)。这句话才是公式的几何含义：剩余误差已经垂直于模型能够调整的两个方向，再沿截距或斜率方向移动都不能继续降低平方误差。

\begin{center}
\begin{tikzpicture}[x=2cm,y=1.25cm,>=stealth]
  \draw[->] (-0.2,0) -- (2.55,0) node[right] {\(x\)};
  \draw[->] (0,-0.1) -- (0,2.65) node[above] {\(y\)};
  \draw[blue!70!black,thick] (-0.1,1.1167) -- (2.35,2.3417)
    node[right] {\(\hat y=7/6+x/2\)};
  \foreach \x/\y/\yh in {0/1/1.1667,1/2/1.6667,2/2/2.1667}{
    \draw[dashed,gray!80] (\x,\y) -- (\x,\yh);
    \fill (\x,\y) circle (2.2pt);
  }
  \node[below] at (1,0) {\(1\)};
  \node[below] at (2,0) {\(2\)};
  \node[left] at (0,1) {\(1\)};
  \node[left] at (0,2) {\(2\)};
  \node[right] at (1.05,1.84) {残差};
\end{tikzpicture}

\emph{最小二乘允许每个点留下残差；它要求的是所有残差合起来不能再被模型方向解释。}
\end{center}

但这行公式仍然藏着条件：\(X\) 的列是否独立，形成 \(X^{\trans}X\) 会不会放大条件数，数据是否足以支持``这条线代表真实规律''。代数上算出一条线、数值上稳定地算出它、统计上相信它能够外推，是三种不同的判断。Beyond Formulas 要做的，就是把这些被一行公式压在一起的问题重新分开。

Beyond Formulas 想做的事其实很朴素：把这一行拆开，看清每一层折叠进去的结构。它不是绕开公式，而是让公式恢复本来的厚度——它压缩了哪些对象和关系，依赖什么条件，经过了怎样的推理，能不能被可靠地计算，又需要什么证据才能支撑一个判断。

全书反复追问的只有六件事。你可以把它们想象成六个镜头，从不同距离对准同一个数学陈述：

\begin{enumerate}
\tightlist
\item
  \textbf{对象是什么？}——变量、函数、分布、空间、数据，各自属于什么类型。连类型都说不清，后面没法谈。
\item
  \textbf{结论为什么成立？}——定义、假设和推导怎样把前提拧到结论上。这里没有魔法，只有必须逐行检查的逻辑链。
\item
  \textbf{结果能算出来吗？}——算法收敛吗？误差会被放大还是缩小？换了一台机器、换了一种精度，答案还一样吗？
\item
  \textbf{证据到底支持什么？}——有限样本允许我们说什么，又禁止我们说什么。不要拿相关当因果，也不要拿点估计当确定性。
\item
  \textbf{边界在哪儿？}——条件变一下、分布偏移一点、系统进入反馈回路之后，方法还站得住吗？知道什么时候会失效，比知道什么时候管用更需要勇气。
\item
  \textbf{结构是什么？}——这个现象背后的低维形状、对称性、不变量是什么？为什么世界总是倾向于用这么少的结构表达那么复杂的行为？
\end{enumerate}

前五个问题管的是``数学怎样撑住计算和判断''——它们让你做得对。第六个问题管的是``数学揭示了世界的什么''——它让你理解自己到底在做一件什么事。两者缺一个都会走偏：只会算而不理解，等于蒙着眼走路；只理解而不会算，等于知道远方有路却迈不出脚。

所以这里不把数学基础当成学完就可以丢掉的预备课。数学语言、空间结构、不确定性、数值计算、优化、学习、控制——它们不会只出现一次。它们会在同一个实际问题的不同阶段反复撞上你，后面的东西会不断把你推回前面，逼你重新审视你以为已经吃透的一个量词、一种收敛、一次矩阵分解。

这套笔记也不声称能穷尽你未来需要的每一个主题。它给的是一个可以扩展的骨架。以后撞上新知识时，你知道该问：它解决的是哪一类困难？它踩在哪些已有结构上面？它应该接到哪条主线上？它和旁边的方案有什么区别？笔记的边界可以往外长，但判断知识位置的坐标系应该是稳定的。

接下来的 \hyperref[ch:002]{``全书知识地图：从对象到决策''} 会展开这个坐标系；\hyperref[ch:003]{``怎样选择学习路线''} 给了几种走法。如果你手头已经有一个让你卡住的具体问题，可以直接跳进 \hyperref[ch:138]{``按问题查找：从一次困惑进入数学''}；如果你更想知道投影、对偶、不动点或不变量这些结构怎么横跨学科，那就从 \hyperref[ch:139]{``按洞见查找：同一种结构怎样跨越学科''} 开始。

书的最后还留了一个部分：\hyperref[ch:143]{``收束：每一章都懂了之后，它们合起来在回答什么''}。它不讲新数学，只把前面所有内容换两种方式重排一遍——一次按一个系统从诞生到上线的时间顺序，一次按那几个反复出现的数学动作。如果你读到某个阶段觉得``每一章都懂了，但不知道它们合起来在干什么''，那一部分就是为这一刻准备的，不必等到全书读完再翻。
